% 中文
\cabstract{
根据实验内容 \cite{lab02}, 在此次实验需要实现一个议程管理系统，此报告参照评分标准\cite{labeval02}，完成这份实验报告撰写并附上相应的实验报告说明。
本实验实现了一个基于命令行（CLI）的议程管理系统（Agenda），使用 Java 语言并采用面向对象设计。系统提供 `register`、`add`、`query`、`delete`、`clear`、`batch`、`quit` 等命令，支持交互式与批处理模式。核心采用单例的 `AgendaManager` 负责命令解析、鉴权与业务调度，命令通过 `Command` 接口实现分发；会议信息以内存结构（`Map<String, User>` 与 `List<Meeting>`）存储，时间处理基于 `java.time`。实现要点包括时间格式校验与解析、时间重叠检测、标签唯一性校验与批处理逐行回放。为保证质量，编写了 JUnit5 单元测试并集成到构建脚本中。本实验强化了面向对象建模、错误处理与自动化测试的实践，为后续引入持久化与并发支持提供了基础。
}
\ckeywords{议程管理；Java；命令行；面向对象；单元测试；批处理}

